\section{第 7 章 总结与展望}

\subsection{7.1 全文工作总结}

本文围绕“基于 Flask 的智能气象信息服务系统设计与实现”这一主题，完成了一套面向杭州地区的综合气象信息服务平台设计与开发。与传统仅提供单一天气展示的页面不同，本系统将多模式预报分析作为主线，把探空诊断、历史气候统计、机器学习误差校正和农业气象服务组织到同一平台之中，形成了从未来预报展示、专业环境分析、历史背景支撑到行业应用延伸的较完整业务闭环。

本文使用Flask构建统一的Web应用框架，用模板页面和接口相结合的方式来组织前后端交互。系统按照路由层、服务层、分析层和模型层对系统模块进行拆分，使得多模式预报、探空分析、历史气候、农业服务等各个功能在统一入口下可以协同工作，又可以保持自己独立的数据处理、分析逻辑。该种实现方式既满足了毕业设计系统完整性要求，又为以后继续增加站点、模式、行业等留下余地。

系统对Open-Meteo接入了ECMWF、GFS、ICON等模式数据，可以做72小时的精细化预报以及多个模式的对比。用户可以查看未来逐小时气象要素的演变，也可以用多模式并列结果来观察各个预报方案之间的一致性以及分歧程度，从而提高对天气过程强弱和转折时段的判断能力。本系统的这个主要模块又引入了 ECMWF 的机器学习误差校正链路，使得预报分析从原来的仅仅是对原始模式结果的展示变成了有一定的后处理优化能力。

系统可以对怀俄明大学探空资料进行层结解析、稳定度参数计算、专业图形展示。通过 Matplotlib 与 MetPy 生成的 Skew-T、T-lnP 等图形，系统能够从垂直方向刻画大气温湿风结构，并通过 CAPE、CIN、K 指数、Lifted Index、可降水量和风切变等参数辅助判断对流潜势和稳定度环境。该模块使系统具备了比常规天气页面更强的诊断能力，也使平台可以为强对流天气分析、农业防灾、航空气象保障等更专业的应用提供支持。

历史气候分析能力系统完成年度页、趋势页、代表性极端事件识别逻辑的实现。本文除了保留原始极值榜之外，又加入了基于历史同月同日经验百分位的EW指数，并且用季节门槛、过程去重和四季分组的方式呈现，使得年度极端事件筛选的结果比较符合统计解释性以及业务可读性。这样历史气候模块就不再是简单的展示几张气候图表，而是可以给当前年份的异常天气过程提供更有力的历史背景支持。

本文选取 ECMWF/Open-Meteo 预报链路作为研究对象，在此基础上对杭州国家站/馒头山站点口径重新组织训练样本，统一了训练期和运行期的站点配置、issue time 识别规则以及评估划分方式。新模型对于温度、湿度、风速等连续变量的误差改善比较明显，在风速订正方面改善幅度较大。虽然降水分类指标并没有在所有的维度上都比原始预报好，但是降水量级的误差明显减少，说明该模块有比较好的业务应用价值。更重要的是，该部分工作使系统预报后的处理过程更加真实、统一、可追溯。

在农业气象服务上，系统把天气预报、历史背景、作物知识库三者融合起来，创建起服务于农业生产场景的模块。系统可以针对水稻、茶树、杨梅、柑橘、小白菜等主要农作物给出积温进度预估、适宜性评定、灾害预警及农事建议的输出。历史积温用ERA5-Land再分析资料估算，未来短期积温增量用ECMWF预报链路推算，使农业模块既能考虑长期背景又能考虑短期决策。由此可知，本系统不但具有综合天气分析的能力，而且也有望向具体行业服务转化。

综合来看，本文完成的主要工作并不在于提出全新的算法，而在于围绕真实业务场景，把多源气象数据、统计分析方法、图形可视化和误差校正能力组织成一套可运行、可展示、可解释的综合系统。对于毕业设计而言，这种“以系统实现为主、以分析能力为核心、以业务闭环为目标”的完成方式具有较强的实践意义。

\subsection{7.2 当前系统的不足}

虽然本文已经完成较为完整的系统实现，但是从实际运行情况来看，目前系统还存在着一些不足，主要表现在数据源稳定性、样本数量、结果泛化能力以及部分模块工程鲁棒性等各方面。

第一对外部的数据源还存在很强的依赖性。不论是Open-Meteo预报接口、怀俄明大学探空服务还是历史观测和再分析资料接口，它们的可用性以及响应稳定性都不是由本地系统所决定的。当外部接口出现访问异常、字段变更或者时次缺少的时候，尽管系统可以依靠缓存、回退以及异常提示来缩减影响，但是依然不可避免地会对结果的完备性以及实时性产生干扰。因此目前的系统在“平台能力自主可控”上还存在不足。

其次，机器学习误差校正模块所用的数据样本量太小。虽然本次用新版杭州站ECMWF预报样本进行统一口径重训，训练和测试的程序也比之前版本要更严格一些，但是整体样本时段还比较短，还不足以覆盖各个季节、各种天气类型以及较长的时长下复杂的误差特征。降水变量本身随机性强、局地性强，在有限样本下很难同时达到分类表现和量级回归的表现，所以目前模型更应该被看作是“初步可用的业务增强工具”，而不是成熟的普适订正方案。

再次，系统目前仍然以杭州单站以及周边区域为场景，泛化能力较差。不论是历史气候分析的口径、机器学习误差校正模型，还是农业气象模块中作物知识的配置，都是以杭州地区为研究对象。聚焦可以提高系统的实现完整性，但是也意味着如果把系统直接迁移到其他地区，很多站点参数、阈值规则、作物结构和训练样本都要重新适应。因此，目前版本更适合用作“区域化综合气象服务系统原型”，不能直接推广到全国任意地区。

另外一些分析逻辑虽然已经满足了目前业务的需求，但是仍然存在一定的简化性。历史气候模块中的EW指数仍然以经验百分位、规则约束为基础，探空模块中部分指数在异常情况下还要依靠回退、空态处理，农业模块中的适宜度、风险阈值也大多依据知识规则、经验来设定。这些做法对于毕业设计系统是合理的，但是要使业务化程度更高，还需要在数据驱动建模上更加细致，在业务验证上更加严格。

系统对于文档化、长期维护的支持性还比较弱。目前项目已经形成了一套比较完整的页面、接口、测试、模型成果，但是对于部署说明、数据来源说明、异常场景处理清单以及后续维护规范，还没有形成一个系统化的整理。对于毕业设计来说，该问题并不会影响主要成果的展示，但是从长远的工程演进来看，文档化不足会增大后续的迭代成本。

\subsection{7.3 后续展望}

针对以上不足，本文认为本系统以后可以从数据、模型、功能、应用四个方面继续前进。

第一，从数据层面加大数据源的稳定性、冗余性。对数值预报、探空、观测数据可以创建比较明确的本地缓存、归档系统，逐渐降低对实时外部接口单点可用性的依靠，对于重要的站点可以增加一些权威的、长时间的数据作为本地观测资料，来提高历史分析以及机器学习训练的基础质量。如果将来条件允许的话，可以逐渐形成统一的数据清单和元数据管理机制，使系统数据来源、时段覆盖、使用边界更清楚。

第二，在模型上可以继续扩大机器学习误差校正的样本量和建模深度。随着越来越多的ECMWF历史预报文件以及更长时间的观测真值被收集起来，之后就可以重新训练更加稳定的订正模型，并且按照季节、按时效或者按照天气类型来进行分组建模，从而更好地适应江南地区中远期预报中经常出现的偏暖、偏湿等系统性误差。除了随机森林之外，还可以研究梯度提升树等更适合表格数据的模型方案，但是前提必须是保持站点口径、时次规则的一致性，不能因为算法复杂度的提高而降低结果的可信度。

第三，系统可以进一步发展成为功能更加专业的、更加完备的分析平台。在多模式预报上可以加入更多的集合信息展示以及模式一致性诊断，在探空上可以补充更多的参数产品、典型天气型识别和图文联动的解释，在历史气候上可以丰富趋势检验、分布特征和极端事件对比的能力，在农业上可以增加更多的作物种类、病虫害气象风险和分阶段农事管理模板。由此系统可以由原来的“综合查询与分析平台”发展成现在的“专题化服务平台”。

第四，从应用角度来讲，系统以后会加强智能分析和自动化服务的功能。目前农业模块已经采用AI建议和规则回退并行的方式，在此基础上，也可以把这种思路应用到多模式预报摘要、探空诊断提示、历史气候年度报告自动生成等场景中，使得系统不仅能给出原始的图表、参数的结果，还能给出更加符合用户决策语言的解释性文字。如果再加上自动化报表和定时任务的能力，那么系统就更有可能成为一个面向农业管理、地方气象科普或者专题业务分析的轻量级服务平台。

本文所完成的系统证明了以多模式预报分析为主、以历史气候和探空诊断为支撑、以机器学习后处理和农业服务为延伸的平台建设思路是可行的。未来工作重点不能只是堆砌新的功能，而应该从数据质量、结果可信度、区域适配性、场景化服务深度四个方面进行持续改进。只有在各个方面不断完善之后，系统才可能由毕业设计原型逐渐发展成一个更加稳定、实用的智能气象信息服务平台。

